%%%%%%%%%%%%%%%%%%%%%%%%%%%%%%%%%%%%%%%%%%%%%%%%%%%%%%%%%%%%%%%%%
%%%
%%% Praca magisterska: Model systemu Smart Home
%%% Autor: Igor Guła
%%% Wydzial Elektrotechniki i Informatyki, Politechnika Rzeszowska
%%%
%%%%%%%%%%%%%%%%%%%%%%%%%%%%%%%%%%%%%%%%%%%%%%%%%%%%%%%%%%%%%%%%%

\documentclass[12pt,twoside]{article}

\usepackage{weiiszablon}

\author{Igor Guła}
\studentID{169784}

\title{Model systemu Smart Home}
\titleEN{Smart Home System Model}

% 2 = praca magisterska
\newcommand{\rodzajPracyNo}{2}

\supervisor{dr inż. Michał Markiewicz}
\supervisorEN{Michał Markiewicz, PhD}

\abstract{
Praca prezentuje projekt oraz implementację platformy SymulatorSH — modelu systemu smart home
umożliwiającego symulację zużycia energii elektrycznej w gospodarstwach domowych oraz analizę
opłacalności różnych taryf energetycznych. System składa się z modułowego symulatora dwunastu
typów urządzeń AGD/RTV, backendu w architekturze REST (Spring Boot 3.3 / Java 21), interfejsu
webowego (React 18 / TypeScript) oraz publicznego API zgodnego ze standardem OpenAPI 3.0.
Aplikacja została wdrożona w chmurze AWS EC2 i integruje się z bazami PostgreSQL (Neon Cloud),
InfluxDB Cloud oraz brokerem MQTT. Jako studium zastosowania modelu przeprowadzono analizę
opłacalności dynamicznej taryfy RDN (Rynek Dnia Następnego) względem tradycyjnych taryf
G11 i G12, wykorzystując rzeczywiste ceny hurtowe z Polskich Sieci Elektroenergetycznych
dla sześciu profili gospodarstw domowych w czterech sezonach meteorologicznych.
}

\abstractEN{
This thesis presents the design and implementation of the SymulatorSH platform — a smart
home system model enabling simulation of household electricity consumption and analysis of
different energy tariff profitability. The system consists of a modular simulator of twelve
types of household appliances, a REST backend (Spring Boot 3.3 / Java 21), a web interface
(React 18 / TypeScript), and a public API compliant with OpenAPI 3.0 standard. The application
has been deployed to AWS EC2 cloud and integrates with PostgreSQL (Neon Cloud), InfluxDB
Cloud databases, and an MQTT broker. As a case study, an analysis of dynamic RDN (Day-Ahead
Market) tariff profitability was performed against traditional G11 and G12 tariffs, using
real wholesale prices from Polish Power Grid for six household profiles across four
meteorological seasons.
}

\keywords{smart home, symulacja, dynamiczne taryfy energii, REST API, mikrousługi}
\keywordsEN{smart home, simulation, dynamic energy tariffs, REST API, microservices}


\begin{document}

% strona tytulowa
\maketitle

\blankpage

% spis tresci
\tableofcontents

\clearpage
\blankpage


\section*{Wykaz skrótów i akronimów}
\addcontentsline{toc}{section}{Wykaz skrótów i akronimów}

\begin{tabular}{ll}
API      & Application Programming Interface \\
AWS      & Amazon Web Services \\
CVaR     & Conditional Value at Risk (warunkowa wartość zagrożona) \\
G11      & Taryfa jednostrefowa (stały cennik) \\
G12      & Taryfa dwustrefowa (dzień/noc) \\
HTTP     & Hypertext Transfer Protocol \\
JSON     & JavaScript Object Notation \\
JWT      & JSON Web Token \\
MQTT     & Message Queuing Telemetry Transport \\
PSE      & Polskie Sieci Elektroenergetyczne \\
RDN      & Rynek Dnia Następnego (dynamiczna taryfa energii) \\
REST     & Representational State Transfer \\
SSE      & Server-Sent Events \\
URE      & Urząd Regulacji Energetyki \\
VaR      & Value at Risk (wartość zagrożona) \\
\end{tabular}


%%%%%%%%%%%%%%%%%%%%%%%%%%%%%%%%%%%%%%%%%%%%%%%%%%%%%%%%%%%%%%%%%
%  ROZDZIAL 1: WSTEP
%%%%%%%%%%%%%%%%%%%%%%%%%%%%%%%%%%%%%%%%%%%%%%%%%%%%%%%%%%%%%%%%%
\section{Wstęp}

TODO: 2 strony — zagajenie, cel pracy, zakres zadań, struktura pracy.


%%%%%%%%%%%%%%%%%%%%%%%%%%%%%%%%%%%%%%%%%%%%%%%%%%%%%%%%%%%%%%%%%
%  ROZDZIAL 2: PODSTAWY TEORETYCZNE
%%%%%%%%%%%%%%%%%%%%%%%%%%%%%%%%%%%%%%%%%%%%%%%%%%%%%%%%%%%%%%%%%
\section{Podstawy teoretyczne}

TODO: 5 stron — rynek energii w Polsce, technologie użyte, przegląd literatury.

\subsection{Rynek energii elektrycznej w Polsce}
\subsection{Wykorzystane technologie}
\subsection{Przegląd literatury}


%%%%%%%%%%%%%%%%%%%%%%%%%%%%%%%%%%%%%%%%%%%%%%%%%%%%%%%%%%%%%%%%%
%  ROZDZIAL 3: METODOLOGIA OBLICZEN
%%%%%%%%%%%%%%%%%%%%%%%%%%%%%%%%%%%%%%%%%%%%%%%%%%%%%%%%%%%%%%%%%
\section{Metodologia obliczeń}

TODO: 5-7 stron ze wzorami — modele kosztów, VaR/CVaR, analiza sezonowa.

\subsection{Model kosztu taryfy G11}
\subsection{Model kosztu taryfy G12}
\subsection{Model kosztu taryfy RDN}
\subsection{Statystyki ryzyka}
\subsection{Metodologia analizy sezonowej}


%%%%%%%%%%%%%%%%%%%%%%%%%%%%%%%%%%%%%%%%%%%%%%%%%%%%%%%%%%%%%%%%%
%  ROZDZIAL 4: ARCHITEKTURA SYSTEMU
%%%%%%%%%%%%%%%%%%%%%%%%%%%%%%%%%%%%%%%%%%%%%%%%%%%%%%%%%%%%%%%%%
\section{Architektura systemu SymulatorSH}
\label{sec:architektura}

TODO: 5-6 stron — diagram komponentów, model danych, API.

\subsection{Diagram komponentów}
\subsection{Stos technologiczny}
\subsection{Model danych}
\subsection{Publiczne REST API}


%%%%%%%%%%%%%%%%%%%%%%%%%%%%%%%%%%%%%%%%%%%%%%%%%%%%%%%%%%%%%%%%%
%  ROZDZIAL 5: IMPLEMENTACJA
%%%%%%%%%%%%%%%%%%%%%%%%%%%%%%%%%%%%%%%%%%%%%%%%%%%%%%%%%%%%%%%%%
\section{Implementacja aplikacji}
\label{sec:implementacja}

Niniejszy rozdział prezentuje warstwę interakcji użytkownika z~systemem SymulatorSH. W~przeciwieństwie do rozdziału~\ref{sec:architektura}, który skupiał się na wewnętrznej strukturze komponentów, poniższe podsekcje opisują aplikację od strony osoby korzystającej z~interfejsu webowego. Każda podsekcja odpowiada jednemu ekranowi lub grupie ekranów aplikacji i~ilustruje realizację konkretnego elementu procesu badawczego: przeglądania wyników, definiowania konfiguracji testu, uruchamiania serii pomiarów oraz analizy kosztów w~trzech taryfach.

\subsection{Widok listy testów}
\label{sec:impl:lista}

Ekran głównym aplikacji stanowi lista testów pogrupowana według partii oraz daty utworzenia (rysunek~\ref{fig:lista_testow}). Grupowanie zostało dodane w~odpowiedzi na obserwację, że po kilkudziesięciu uruchomieniach płaska lista traci czytelność. Nazwa grupy odzwierciedla źródło uruchomienia: prefiksy w~nawiasach kwadratowych (np.~\texttt{[Partia 2026-08-06T18:57]}) oznaczają testy uruchomione zbiorczo z~poziomu ekranu partii, prefiks \texttt{[example.py]} identyfikuje testy uruchomione za pomocą skryptu klienckiego korzystającego z~publicznego API, natomiast grupy datowe zawierają testy uruchomione pojedynczo z~poziomu kreatora lub trybu JSON.

\begin{figure}[htbp]
    \centering
    \includegraphics[width=0.9\textwidth]{figures/lista_testow.png}
    \caption{Lista testów pogrupowana według partii i~daty utworzenia.}
    \label{fig:lista_testow}
\end{figure}

Każdy wiersz w~rozwiniętej grupie zawiera nazwę testu, jego status, liczbę dni symulowanych, współczynnik przyspieszenia oraz znaczniki czasu utworzenia i~zakończenia. Kolorowe etykiety statusu (Zakończony, Błąd) przyspieszają lokalizację nieudanych uruchomień. Odświeżanie listy realizowane jest co pięć sekund, co pozwala obserwować postęp równolegle wykonywanych testów bez konieczności ręcznego przeładowania strony. Kliknięcie w~identyfikator testu otwiera ekran szczegółów opisany w~sekcji~\ref{sec:impl:szczegoly}.

\subsection{Kreator konfiguracji i tryb JSON}
\label{sec:impl:kreator}

Zdefiniowanie konfiguracji testu wymaga wskazania listy urządzeń wraz z~ich parametrami oraz harmonogramem pracy. Aplikacja udostępnia dwie równolegle dostępne ścieżki: kreator formularzowy oraz tryb edycji surowego dokumentu JSON.

Kreator (rysunek~\ref{fig:kreator}) prowadzi użytkownika przez wypełnienie kolejnych sekcji: parametry ogólne (nazwa testu, opis, długość symulacji, współczynnik przyspieszenia, poziomy jittera czasowego i~mocowego) oraz listę urządzeń. Każde urządzenie posiada określony typ (oświetlenie, chłodziarka, czajnik i~inne), przypisane pomieszczenie oraz identyfikator. Rozwijane pola parametrów zawierają wartości domyślne zaczerpnięte z~profilowych ustawień, a~edytor harmonogramu pozwala definiować dowolną liczbę przedziałów aktywności w~formacie \texttt{od-do}. Górny pasek udostępnia mechanizm szablonów, umożliwiający zapisanie wypełnionego formularza jako gotowej konfiguracji do wielokrotnego użycia oraz późniejsze wczytanie jej jednym kliknięciem.

\begin{figure}[htbp]
    \centering
    \includegraphics[width=0.9\textwidth]{figures/kreator_testu.png}
    \caption{Kreator testu z~zaznaczonym szablonem profilu B (Pracownik zdalny) w~sezonie zimowym.}
    \label{fig:kreator}
\end{figure}

Dla zaawansowanych użytkowników oraz na potrzeby integracji zewnętrznej udostępniono tryb surowy (rysunek~\ref{fig:json}). Prezentuje on identyczną strukturę danych, jaką backend akceptuje w~zapytaniach do publicznego API, dzięki czemu ta sama konfiguracja może zostać przygotowana w~przeglądarce, a~następnie skopiowana do skryptu klienckiego (lub odwrotnie). Rozwiązanie to okazało się przydatne podczas testów rekonstruowanych z~logów oraz przy przygotowywaniu przykładów kodu do dokumentacji.

\begin{figure}[htbp]
    \centering
    \includegraphics[width=0.9\textwidth]{figures/JSON.png}
    \caption{Tryb JSON umożliwiający ręczną edycję konfiguracji zgodnej ze specyfikacją publicznego API.}
    \label{fig:json}
\end{figure}

Oba tryby generują dokładnie tę samą strukturę wejściową, walidowaną po stronie backendu. W~razie błędu walidacji użytkownik otrzymuje komunikat wskazujący brakujące lub niepoprawne pola przed uruchomieniem testu.

\subsection{Uruchamianie partii testów}
\label{sec:impl:partia}

Aby usprawnić generowanie danych do analizy porównawczej, wprowadzono ekran zbiorczego uruchamiania partii (rysunek~\ref{fig:partia}). Umożliwia on zaznaczenie dowolnego podzbioru z~macierzy sześciu profili gospodarstw domowych i~czterech sezonów, uzupełnionego o~zapisane szablony własne. Dla każdej pary profil--sezon backend tworzy niezależny test, przypisując mu wspólny identyfikator partii ułatwiający późniejsze zbiorcze pobranie wyników.

\begin{figure}[htbp]
    \centering
    \includegraphics[width=0.9\textwidth]{figures/partia_testow.png}
    \caption{Ekran partii testów z~zaznaczonymi profilami B i~C we wszystkich sezonach oraz wybranymi szablonami własnymi.}
    \label{fig:partia}
\end{figure}

Panel parametrów pozwala ustawić wspólną długość symulacji (od jednego dnia dla testu dymnego do trzydziestu dni dla pełnego badania), współczynnik przyspieszenia oraz interwał emisji pomiarów. Ostatnia z~wymienionych opcji ma bezpośredni wpływ na liczbę punktów zapisywanych do bazy szeregów czasowych, a~pośrednio~-- na obciążenie limitów planu bezpłatnego InfluxDB Cloud. Sekcja podsumowująca w~stopce prezentuje trzy informacje operacyjne: liczbę wybranych testów, szacowany czas rzeczywisty wykonania (uwzględniający wielkość puli wykonawców) oraz przewidywaną liczbę punktów pomiarowych. Wartości te wyliczane są w~przeglądarce na bieżąco, w~oparciu o~parametry konfiguracji, co pozwala świadomie skalować rozmiar partii przed zatwierdzeniem uruchomienia.

Backend wykonuje testy równolegle w~puli czterech wątków. Pozostałe zadania oczekują w~kolejce i~są uruchamiane w~miarę zwalniania zasobów. Podejście to zostało dobrane empirycznie jako kompromis między czasem realizacji partii a~stabilnością zapisu do bazy szeregów czasowych na współdzielonej instancji chmurowej.

\subsection{Ekran szczegółów testu}
\label{sec:impl:szczegoly}

Ekran szczegółów pojedynczego testu prezentuje pełny zestaw metadanych oraz interaktywne wizualizacje pomiarów zgromadzonych w~bazie szeregów czasowych. Nagłówek zawiera nazwę testu, opis pochodzący z~konfiguracji, cztery panele podsumowujące (status, liczba dni symulacji, współczynnik przyspieszenia, sumaryczne zużycie energii) oraz znaczniki czasu utworzenia i~zakończenia, a~także rzeczywisty czas wykonania w~sekundach.

Pod nagłówkiem umieszczono komponent wizualizacji poboru mocy w~czasie. Domyślnie prezentuje on wykres zbiorczy wszystkich urządzeń (rysunek~\ref{fig:profil_razem}), pozwalający ocenić profil zużycia całego gospodarstwa i~zlokalizować dobowe szczyty. Kolory poszczególnych urządzeń są spójne z~legendą u~dołu wykresu, a~oś czasu może być prezentowana jako czas symulowany lub czas rzeczywisty.

\begin{figure}[htbp]
    \centering
    \includegraphics[width=0.9\textwidth]{figures/profil_F_wykres_1.png}
    \caption{Zbiorczy wykres poboru mocy dla profilu F (Para bez dzieci) w~sezonie jesiennym --- 30~dni symulacji, 77~989 pomiarów.}
    \label{fig:profil_razem}
\end{figure}

Przełącznik \emph{Osobno} rozdziela wykres na siatkę wykresów jednostkowych pogrupowanych według pomieszczeń (rysunek~\ref{fig:profil_osobno}). Widok ten umożliwia jakościową ocenę pracy pojedynczego urządzenia: widoczne są cykliczne włączenia zmywarki i~piekarnika, krótkie impulsy czajnika oraz stałe tło poboru mocy chłodziarki. Dodatkowy przełącznik skali (liniowa lub logarytmiczna) pomaga porównywać urządzenia o~znacząco różnym rzędzie wielkości poboru mocy w~jednym widoku.

\begin{figure}[htbp]
    \centering
    \includegraphics[width=0.9\textwidth]{figures/profil_F_wiele_wykresow.png}
    \caption{Wykres poboru mocy w~trybie \emph{Osobno} --- pięć urządzeń przypisanych do pomieszczenia \emph{Kuchnia}.}
    \label{fig:profil_osobno}
\end{figure}

Uzupełnieniem wykresu poboru mocy jest widok \emph{Wzorzec dnia} (rysunek~\ref{fig:wzorzec}), prezentujący uśredniony rozkład aktywności urządzeń w~ciągu doby. Każdy wiersz macierzy odpowiada jednemu urządzeniu, każda kolumna jednej godzinie doby, a~intensywność koloru komórki odzwierciedla średnią moc obserwowaną w~tej godzinie na przestrzeni całego testu (jasny odcień oznacza wyłączenie, pełne nasycenie -- pracę z~mocą znamionową). Widok ten szybko ujawnia charakter danego profilu: ciągłe działanie chłodziarki i~routera, poranno-wieczorną aktywność oświetlenia oraz koncentrację pracy zmywarki i~piekarnika w~godzinach wieczornych. Uśrednianie po wszystkich dobach testu redukuje szum wprowadzony jitterem, przez co wzorzec zachowania jest łatwiejszy do interpretacji niż z~wykresu chwilowego poboru mocy.

\begin{figure}[htbp]
    \centering
    \includegraphics[width=0.9\textwidth]{figures/wzorzec_dnia.png}
    \caption{Wzorzec dobowy urządzeń dla profilu F (Para bez dzieci) --- macierz uśrednionej mocy pogrupowana według pomieszczeń.}
    \label{fig:wzorzec}
\end{figure}

\subsection{Podgląd cen Rynku Dnia Następnego}
\label{sec:impl:ceny}

Odrębny ekran udostępnia interaktywny podgląd godzinowych cen Rynku Dnia Następnego (rysunek~\ref{fig:ceny_rdn}). Po wybraniu daty aplikacja pobiera dobowy komplet dwudziestu czterech notowań z~lokalnej bazy relacyjnej, do której backend cyklicznie replikuje dane z~publicznego API operatora sieci przesyłowej.

\begin{figure}[htbp]
    \centering
    \includegraphics[width=0.9\textwidth]{figures/ceny_energi_rdn.png}
    \caption{Ekran cen RDN dla dnia 16.08.2026 --- średnia dobowa 561~zł/MWh, minimum 91~zł/MWh o~13:00, maksimum 861~zł/MWh o~21:00.}
    \label{fig:ceny_rdn}
\end{figure}

Cztery panele podsumowujące prezentują średnią, wartość minimalną (dolinę cenową), wartość maksymalną (szczyt) oraz stosunek maksimum do minimum. Ostatnia z~wymienionych statystyk jest szczególnie istotna z~perspektywy analizy zmienności: dla prezentowanego dnia stosunek wynosi $\times$9,4, co ilustruje skalę fluktuacji, jakim podlegają ceny w~ciągu doby. Histogram godzinowy pod panelami koduje kolorem trzy kategorie cen wyznaczone względem średniej dobowej: kolor zielony oznacza cenę o~co najmniej 25\% poniżej średniej (okres taniej energii), pomarańczowy~-- cenę zbliżoną do średniej, natomiast czerwony~-- cenę o~co najmniej 25\% powyżej średniej (okres drogiej energii). Uproszczona kategoryzacja pozwala szybko zidentyfikować przedziały czasowe atrakcyjne z~perspektywy odbiorcy na taryfie dynamicznej bez konieczności analizy dokładnych wartości liczbowych.

\subsection{Kalkulator porównawczy taryf}
\label{sec:impl:kalkulator}

Ostatnim komponentem interfejsu jest kalkulator porównawczy, dostępny jako zakładka na ekranie szczegółów testu. Jego zadaniem jest przeliczenie zarejestrowanego przebiegu poboru mocy na koszt energii w~trzech taryfach: G11 (stały cennik), G12 (dwustrefowa dzień/noc) oraz RDN (dynamiczna taryfa godzinowa oparta o~ceny Rynku Dnia Następnego).

Kalkulator udostępnia dwa tryby pracy. Tryb \emph{Test rzeczywisty} (rysunek~\ref{fig:kalkulator_test}) wykorzystuje ceny RDN z~okresu, w~którym test został uruchomiony, i~pełni funkcję porównania punktowego. Tryb \emph{Projekcja historyczna} (nie prezentowany na osobnym rysunku) umożliwia zastosowanie zarejestrowanego profilu poboru do dowolnego historycznego przedziału dat, co pozwala oszacować, jak zachowałby się koszt użytkowania tego samego gospodarstwa w~innym okresie roku lub w~innym roku kalendarzowym.

\begin{figure}[htbp]
    \centering
    \includegraphics[width=0.9\textwidth]{figures/porownanie_cen_profil_F.png}
    \caption{Kalkulator porównawczy taryf w~trybie \emph{Test rzeczywisty} dla profilu F (Para bez dzieci) --- okres 30~dni, sezon jesienny 2025.}
    \label{fig:kalkulator_test}
\end{figure}

Panele wynikowe prezentują sumaryczny koszt energii w~każdej z~trzech taryf, wyróżniając zieloną etykietą taryfę najtańszą dla analizowanego przypadku. Poniżej umieszczono sekcję statystyk dziennego kosztu w~taryfie dynamicznej, obejmującą wartość minimalną, medianę, wartość zagrożoną VaR na poziomie 5\% (próg 5\% najgorszych dób), warunkową wartość zagrożoną CVaR na poziomie 5\% (średnia z~ogona rozkładu) oraz wartość maksymalną (najgorszy scenariusz). Dobór tych statystyk został podyktowany specyfiką taryfy RDN: średnia arytmetyczna nie oddaje w~pełni ryzyka, ponieważ ukrywa incydenty o~ekstremalnie wysokim koszcie dobowym, które mogą w~sposób istotny wpływać na budżet gospodarstwa domowego. Zastosowanie VaR i~CVaR znanych z~analiz finansowych pozwala scharakteryzować zarówno typową ekspozycję, jak i~jej ogon rozkładu.

Poniżej statystyk umieszczono krótki wniosek tekstowy porównujący sumaryczny koszt RDN z~kosztem G11, a~następnie dwa wykresy: słupkowy pokazujący łączny koszt w~trzech taryfach oraz liniowy prezentujący dobowe koszty w~całym analizowanym okresie. Ten ostatni ujawnia charakter różnicy między taryfą stałą a~dynamiczną: linie G11 i~G12 są niemal poziome, natomiast krzywa RDN wykazuje wyraźną zmienność z~pojedynczymi szczytami odpowiadającymi dniom o~podwyższonych cenach hurtowych. Wizualizacja ta stanowi punkt wyjścia do dyskusji nad opłacalnością taryfy dynamicznej prowadzonej w~rozdziale~\ref{sec:wyniki}.


%%%%%%%%%%%%%%%%%%%%%%%%%%%%%%%%%%%%%%%%%%%%%%%%%%%%%%%%%%%%%%%%%
%  ROZDZIAL 6: METODYKA BADAN
%%%%%%%%%%%%%%%%%%%%%%%%%%%%%%%%%%%%%%%%%%%%%%%%%%%%%%%%%%%%%%%%%
\section{Metodyka badań}

TODO: 3-4 strony — profile, scenariusze, proces.

\subsection{Definicja profili gospodarstw domowych}
\subsection{Scenariusze sezonowe}
\subsection{Proces uruchamiania testów}
\subsection{Metodologia re-kalkulacji z cenami historycznymi}


%%%%%%%%%%%%%%%%%%%%%%%%%%%%%%%%%%%%%%%%%%%%%%%%%%%%%%%%%%%%%%%%%
%  ROZDZIAL 7: WYNIKI BADAN
%%%%%%%%%%%%%%%%%%%%%%%%%%%%%%%%%%%%%%%%%%%%%%%%%%%%%%%%%%%%%%%%%
\section{Wyniki badań i analiza}
\label{sec:wyniki}

TODO: 6-7 stron — tabela zbiorcza, analiza per profil, per sezon, trade-off, hipotezy.

\subsection{Tabela zbiorcza wyników}
\subsection{Analiza per profil gospodarstwa}
\subsection{Analiza per sezon}
\subsection{Trade-off zysk vs ryzyko}
\subsection{Analiza sensitivity}
\subsection{Weryfikacja hipotez badawczych}


%%%%%%%%%%%%%%%%%%%%%%%%%%%%%%%%%%%%%%%%%%%%%%%%%%%%%%%%%%%%%%%%%
%  ROZDZIAL 8: PODSUMOWANIE
%%%%%%%%%%%%%%%%%%%%%%%%%%%%%%%%%%%%%%%%%%%%%%%%%%%%%%%%%%%%%%%%%
\section{Podsumowanie i wnioski końcowe}

TODO: 1.5 strony — cel zrealizowany, wnioski, kierunki rozwoju, wkład własny.


%%%%%%%%%%%%%%%%%%%%%%%%%%%%%%%%%%%%%%%%%%%%%%%%%%%%%%%%%%%%%%%%%
%  ZALACZNIKI
%%%%%%%%%%%%%%%%%%%%%%%%%%%%%%%%%%%%%%%%%%%%%%%%%%%%%%%%%%%%%%%%%
\section*{Załączniki}
\addcontentsline{toc}{section}{Załączniki}

TODO: opcjonalnie — fragmenty kodu, konfiguracje, dodatkowe wykresy.


\clearpage

\addcontentsline{toc}{section}{Literatura}

\begin{thebibliography}{99}
\bibitem{TODO} TODO: pozycje literatury dodać w trakcie pisania rozdziałów.
\end{thebibliography}

\clearpage

\makesummary

\end{document}
